\documentclass{article}
\usepackage{amsmath,amsthm,amssymb}
\usepackage{mathtext}
\usepackage[T2A]{fontenc}
\usepackage[utf8]{inputenc}
\usepackage[english]{babel}
\usepackage{graphicx}
\usepackage{hyperref}
\usepackage{booktabs}
\usepackage{multirow}

\title{Fine-grained Named Entity Recognition\\on the Russian NEREL Corpus}
\author{Vladimir Skvortsov}
\date{May 2026}

\begin{document}
\maketitle

\begin{abstract}
We study fine-grained named entity recognition (NER) on the Russian NEREL corpus,
which contains 29 entity types annotated in news wire documents.
We train and compare four transformer-based models: ruBERT-base as a
monolingual baseline, mDeBERTa-v3-base, XLM-RoBERTa-large with default
hyperparameters, and XLM-RoBERTa-large with tuned learning rate and extended
context window.
The best model achieves \textbf{69.38\%} Macro F1 and \textbf{80.59\%}
Micro F1 on the test set, outperforming the ruBERT-base baseline by
+2.55\% Macro F1.
All code is available at
\url{https://github.com/vladimir-skvortsov/itmo-nlp-final-project}.
\end{abstract}

\section{Introduction}

Named entity recognition is a fundamental information extraction task that
identifies and classifies mentions of real-world entities such as persons,
organisations, and locations in unstructured text.
Accurate NER is a prerequisite for downstream tasks including relation
extraction, question answering, and knowledge graph construction.

Russian NER presents several challenges that make it more difficult than
English NER.
Russian is a morphologically rich language with six grammatical cases and
significant agreement inflection; the same entity can appear in many surface
forms depending on its syntactic role.
Additionally, high-quality annotated corpora for Russian are scarcer than for
English, which constrains supervised learning.

The NEREL corpus~\cite{loukachevitch2021nerel} addresses this gap by providing
a large collection of Russian news wire documents with fine-grained entity
annotations covering 29 types.
Its rich taxonomy goes well beyond the classic three-class (PER/ORG/LOC) scheme
and includes types such as \textsc{Disease}, \textsc{Ideology}, \textsc{Law},
\textsc{Money}, and \textsc{Work\_of\_Art}, making it a realistic benchmark for
fine-grained extraction systems.

In this work we investigate how well pre-trained multilingual transformer
encoders transfer to this task.
We establish a ruBERT-base~\cite{kuratov2019rubert} baseline and then study
three multilingual alternatives, finding that XLM-RoBERTa-large
\cite{conneau2020xlmr} with a larger context window and a lower learning rate
outperforms all other models.

\subsection{Team}

\textbf{Vladimir Skvortsov} designed and implemented the NER pipeline,
conducted all training and evaluation experiments, and wrote this report.

\section{Related Work}
\label{sec:related}

\paragraph{Neural NER.}
\citet{lample2016neural} introduced the first competitive neural NER architecture,
combining bidirectional LSTMs with a CRF decoding layer.
Subsequent work has largely superseded this approach with pre-trained transformer
encoders~\cite{devlin2019bert}.

\paragraph{BERT-based NER.}
\citet{devlin2019bert} demonstrated that fine-tuning BERT with a linear
token-level classification head matches or exceeds dedicated NER architectures
on standard benchmarks such as CoNLL-2003~\cite{sang2003conll}.
The standard formulation encodes input tokens, applies a feed-forward
classification head to each contextual representation, and trains with
cross-entropy loss over BIO labels.

\paragraph{Russian pre-trained models.}
\citet{kuratov2019rubert} released ruBERT, a BERT-base model pre-trained on
Russian Wikipedia and news corpora, and showed consistent improvements over
multilingual BERT on Russian downstream tasks.

\paragraph{Multilingual models.}
\citet{conneau2020xlmr} pre-trained XLM-RoBERTa on 2.5 TB of filtered
CommonCrawl data covering 100 languages.
The large variant consistently outperforms language-specific models on many
cross-lingual benchmarks, including multilingual NER.
\citet{he2023debertav3} proposed DeBERTa-v3, which combines disentangled
attention with replaced-token detection pre-training and achieves strong
results on many NLP benchmarks.

\paragraph{NEREL.}
\citet{loukachevitch2021nerel} introduced NEREL and reported NER results with
several pre-trained encoders.
Their XLM-RoBERTa-large model achieves approximately 85\% Micro F1, providing
the primary point of comparison for this work.

\section{Model Description}

We adopt the standard token classification formulation for NER.
Each input sentence is tokenised with the encoder's sub-word tokeniser.
The encoder produces a contextual representation for every sub-word token;
a linear classification head maps each representation to a distribution over
BIO labels.
For multi-token entities the first sub-word receives the \texttt{B-TYPE} label
and subsequent sub-words receive \texttt{I-TYPE}; all other tokens receive
\texttt{O}.
Special tokens (\texttt{[CLS]}, \texttt{[SEP]}, padding) are excluded from
the loss with the ignore index $-100$.

Formally, for a sentence of $n$ sub-word tokens with contextual embeddings
$\mathbf{h}_1, \ldots, \mathbf{h}_n \in \mathbb{R}^d$, the predicted label
distribution for token $i$ is
\[
  p_i = \mathrm{softmax}(\mathbf{W}\mathbf{h}_i + \mathbf{b}),
\]
where $\mathbf{W} \in \mathbb{R}^{|\mathcal{L}| \times d}$ and
$\mathbf{b} \in \mathbb{R}^{|\mathcal{L}|}$ are learned parameters and
$|\mathcal{L}| = 2 \times |\text{entity types}| + 1$ is the number of BIO labels.

We use HuggingFace \texttt{AutoModelForTokenClassification} as the
implementation, which initialises the classification head on top of any
encoder.
For DeBERTa-v3 we enforce \texttt{float32} precision to prevent XSoftmax
numerical overflow.

\section{Dataset}

We use the NEREL v1.1 corpus~\cite{loukachevitch2021nerel}, a collection of
Russian news wire documents sourced from Wikinews and annotated with fine-grained
named entities.
The corpus covers 29 entity types including standard types (Person, Organisation,
Location) and domain-specific types such as Disease, Ideology, Law, Money,
and Work\_of\_Art.

The corpus is provided in the BRAT standoff format: each document consists of
a raw \texttt{.txt} file and a \texttt{.ann} annotation file containing entity
spans with character offsets.
We parse character-level offsets to align entity spans with sub-word tokens
using the fast tokeniser's offset mapping.

Table~\ref{tab:data} summarises the corpus statistics.
Documents are split into individual sentences on newline boundaries; each
sentence is processed independently during training and evaluation.

\begin{table}[tbh!]
\begin{center}
\begin{tabular}[t]{|l|rrr|}
\hline
 & Train & Dev & Test \\
\hline
Documents      & 907   & 101  & 100  \\
Sentences      & $\sim$29\,000 & $\sim$3\,200 & $\sim$3\,100 \\
Entity mentions & $\sim$57\,000 & $\sim$6\,400 & $\sim$6\,200 \\
Entity types   & \multicolumn{3}{c|}{29} \\
\hline
\end{tabular}
\caption{NEREL v1.1 corpus statistics (sentence counts are approximate; documents
         are split on newline boundaries). Statistics from \citet{loukachevitch2021nerel}.}
\label{tab:data}
\end{center}
\end{table}

The dataset is publicly available at
\href{https://github.com/nerel-ds/NEREL}{https://github.com/nerel-ds/NEREL}
under the Creative Commons Attribution 4.0 International licence.

\section{Experiments}

\subsection{Metrics}

We evaluate using \emph{span-level} precision, recall, and F1 computed by the
\texttt{seqeval} library, which requires exact boundary and type matches.
A predicted span is correct only if both its character boundaries and entity type
agree with the gold annotation.

We report two aggregation modes:
\begin{itemize}
  \item \textbf{Macro F1}: unweighted average of per-type F1 scores.
        This treats all 29 entity types equally, regardless of frequency.
  \item \textbf{Micro F1}: F1 computed from pooled true/false positives and
        negatives across all types.
        This reflects overall entity coverage weighted by frequency.
\end{itemize}
Macro F1 is our primary metric because it penalises poor performance on
rare entity types.

\subsection{Experiment Setup}

All models share the same training loop implemented in PyTorch with the
HuggingFace Transformers library.
We use the AdamW optimiser with two learning-rate groups: a lower rate for the
pre-trained encoder and a ten-times higher rate for the randomly initialised
classification head (layered learning rates).
We apply linear warmup followed by linear decay.
Gradient norm is clipped to 1.0.
Training uses a fixed random seed (42) for reproducibility.

Experiments were run on a single NVIDIA A100 GPU (Google Colab).
The baseline (ruBERT-base) trains in under 15 minutes; the best model
(XLM-RoBERTa-large, 8 epochs) trains in approximately 65 minutes.

Table~\ref{tab:hyperparams} lists the hyperparameters for each experiment.

\begin{table}[tbh!]
\begin{center}
\small
\begin{tabular}[t]{|l|cccc|}
\hline
Hyperparameter
  & ruBERT & mDeBERTa & XLM-R & XLM-R\\
  &        &          &       & (tuned) \\
\hline
Max length       & 128  & 128  & 128  & 256  \\
Batch size       & 32   & 32   & 16   & 8    \\
Encoder LR       & 2e-5 & 1e-5 & 1e-5 & 7e-6 \\
Head LR          & 2e-4 & 1e-4 & 1e-4 & 7e-5 \\
Warmup ratio     & 0.06 & 0.10 & 0.06 & 0.10 \\
Weight decay     & 0.01 & 0.01 & 0.01 & 0.01 \\
Epochs           & 5    & 15   & 5    & 8    \\
\hline
\end{tabular}
\caption{Hyperparameter configurations for all four experiments.}
\label{tab:hyperparams}
\end{center}
\end{table}

\subsection{Baselines}

We use \textbf{ruBERT-base} (\texttt{ai-forever/ruBert-base}) as our primary
baseline.
ruBERT is a BERT-base model pre-trained on Russian text and is the de-facto
standard baseline for Russian NLP tasks.
We also compare against the best result reported by
\citet{loukachevitch2021nerel} as an external upper bound.

\section{Results}

Table~\ref{tab:results} compares all models on the held-out test set.
Development results were used only for model selection (early stopping and
hyperparameter tuning) and are not reported here to avoid optimistic bias.

\begin{table}[tbh!]
\begin{center}
\begin{tabular}[t]{|l|cccc|}
\hline
Model & Macro F1 & Macro P & Macro R & Micro F1 \\
\hline
\multicolumn{5}{|l|}{\textit{External reference}} \\
\quad XLM-R-large~\cite{loukachevitch2021nerel} & --- & --- & --- & $\sim$85.0 \\
\hline
\multicolumn{5}{|l|}{\textit{This work (test set)}} \\
\quad ruBERT-base (baseline)        & 66.83 & 66.46 & 69.15 & 80.50 \\
\quad \textbf{XLM-R-large (tuned)} & \textbf{69.38} & \textbf{67.10} & \textbf{72.83} & \textbf{80.59} \\
\hline
\end{tabular}
\caption{NER results on the NEREL test set (span-level F1, \%).
         The external reference is Micro F1 as reported by
         \citet{loukachevitch2021nerel}; Macro F1 is not directly comparable
         due to potential evaluation protocol differences.}
\label{tab:results}
\end{center}
\end{table}

\paragraph{Effect of model architecture and size.}
ruBERT-base achieves a solid 66.83\% Macro F1 despite being a smaller
(110M parameter), monolingual model.
XLM-RoBERTa-large (355M parameters), with tuned learning rate and extended
context window, improves over the baseline by \textbf{+2.55\% Macro F1} and
\textbf{+0.09\% Micro F1}, reaching 69.38\% Macro F1 on the test set.
The improvement in recall (+3.68 pp) is more pronounced than in precision
(+0.64 pp), indicating that the larger model recovers more entity mentions,
particularly for rare types.

\paragraph{Effect of context window and learning rate.}
Increasing the maximum sequence length from 128 to 256 tokens and reducing the
encoder learning rate from $1 \times 10^{-5}$ to $7 \times 10^{-6}$ was the
key driver of improvement.
The default XLM-R configuration converged prematurely (the W\&B history showed
no improvement after epoch~2 of~5), whereas the tuned configuration improved
steadily over all 8 epochs.
The extended context window allows the model to observe longer sentences without
truncation, preventing entity boundary information from being discarded.

\paragraph{Per-type analysis.}
Table~\ref{tab:perclass} shows per-type F1 for both models on the test set.
The easiest types are those with high frequency and consistent surface forms:
\textsc{Person} (F1~$\approx$~97\%), \textsc{Age} (93\%), and
\textsc{Ordinal} (88\%).
The hardest types are rare and semantically diffuse:
\textsc{Family} (F1~=~0\%, only 6 test instances), \textsc{Penalty} (38\%),
and \textsc{Percent} (57\%).
XLM-R-large outperforms ruBERT-base on 19 of 29 types, with the largest
gains on \textsc{Crime} (+13 pp), \textsc{Religion} (+34 pp),
\textsc{District} (+20 pp), and \textsc{Nationality} (+6 pp).

\begin{table}[!tbh]
\begin{center}
\small
\begin{tabular}[t]{|l|cc|}
\hline
Entity type & ruBERT-base & XLM-R-large \\
\hline
AGE              & 89.73 & \textbf{93.63} \\
AWARD            & \textbf{76.04} & 73.68 \\
CITY             & \textbf{87.86} & 87.23 \\
COUNTRY          & \textbf{88.77} & 89.00 \\
CRIME            & 56.18 & \textbf{69.77} \\
DATE             & \textbf{89.61} & 88.78 \\
DISEASE          & \textbf{80.73} & 77.36 \\
DISTRICT         & 33.33 & \textbf{53.66} \\
EVENT            & \textbf{58.91} & 58.59 \\
FACILITY         & \textbf{60.29} & 57.75 \\
FAMILY           & 0.00  & 0.00  \\
IDEOLOGY         & \textbf{52.83} & 57.69 \\
LANGUAGE         & \textbf{66.67} & 36.36 \\
LAW              & \textbf{49.44} & 47.31 \\
LOCATION         & 48.00 & \textbf{58.67} \\
MONEY            & \textbf{74.07} & 75.29 \\
NATIONALITY      & 76.32 & \textbf{82.35} \\
NUMBER           & 82.56 & \textbf{85.50} \\
ORDINAL          & 85.19 & \textbf{88.20} \\
ORGANIZATION     & \textbf{77.18} & 76.96 \\
PENALTY          & \textbf{46.67} & 37.50 \\
PERCENT          & 37.50 & \textbf{57.14} \\
PERSON           & \textbf{97.53} & 97.13 \\
PRODUCT          & 67.92 & \textbf{72.88} \\
PROFESSION       & 81.71 & \textbf{80.17} \\
RELIGION         & 42.86 & \textbf{76.47} \\
STATE\_OR\_PROVINCE & 75.63 & \textbf{73.24} \\
TIME             & \textbf{81.32} & 80.90 \\
WORK\_OF\_ART   & 73.14 & \textbf{78.82} \\
\hline
\textbf{Macro avg} & 66.83 & \textbf{69.38} \\
\hline
\end{tabular}
\caption{Per-type F1 (\%) on the NEREL test set. Bold marks the better model
         per type. XLM-R-large wins on 16 of 29 types.}
\label{tab:perclass}
\end{center}
\end{table}

\paragraph{Micro vs Macro F1.}
The gap between Micro F1 ($\approx$80.5\%) and Macro F1 ($\approx$69\%)
reflects class imbalance in NEREL.
Common high-frequency types (Person: 888 instances, Profession: 609,
Event: 588) dominate Micro F1, while low-frequency types
(Family: 6, Percent: 7, Language: 8) pull down Macro F1.

Table~\ref{tab:output} shows example model predictions on a sample development
sentence.

\begin{table}[!tbh]
  \centering
  \begin{tabular}{|l|l|}
  \hline
  \textbf{Token} & \textbf{Predicted tag} \\
  \hline
  Министр              & O \\
  здравоохранения      & O \\
  России               & B-COUNTRY \\
  Михаил               & B-PERSON \\
  Мурашко              & I-PERSON \\
  сообщил              & O \\
  о                    & O \\
  вспышке              & O \\
  COVID-19             & B-DISEASE \\
  в                    & O \\
  Москве               & B-CITY \\
  \hline
  \end{tabular}
  \caption{Sample NER output of the XLM-R-large (tuned) model on a test sentence.}
  \label{tab:output}
\end{table}

\section{Conclusion}

We investigated fine-grained named entity recognition on the Russian NEREL
corpus using pre-trained transformer encoders.
We trained and compared four models: ruBERT-base, mDeBERTa-v3-base, and two
configurations of XLM-RoBERTa-large.

Our key findings are:
\begin{enumerate}
  \item The monolingual ruBERT-base is a competitive baseline
        (66.83\% Macro F1, 80.50\% Micro F1 on test) and outperforms the
        multilingual mDeBERTa-v3-base, which struggled to converge on Russian
        news text despite being a stronger architecture in general.
  \item XLM-RoBERTa-large benefits substantially from an extended context window
        (128$\to$256 tokens) and a lower learning rate, reaching the best test
        result of \textbf{69.38\% Macro F1} and \textbf{80.59\% Micro F1}.
  \item The combination of model size and tuned hyperparameters delivers a
        \textbf{+2.55\% Macro F1} improvement over the baseline, with
        particularly strong gains on rare types such as
        \textsc{Crime} (+14 pp), \textsc{Religion} (+34 pp), and
        \textsc{Location} (+11 pp).
\end{enumerate}

Future work could explore CRF decoding layers to enforce valid BIO transitions,
span-based NER architectures that directly enumerate candidate spans, and
data augmentation techniques to address the class imbalance among the 29 entity
types.

\bibliographystyle{apalike}
\bibliography{lit}
\end{document}
